\section{Conclusion}
\label{sec:conclusion}

Bounded addition and inversion can be implemented by a sparse resistive
network that retains every source solution and every coordinate needed to
recover it. As a result, exact resistive feasibility is $\ER$-complete even
under the simultaneous restrictions of Theorem~\ref{thm:structural}, in a
model with independently prescribed voltage and injection bounds, including
simultaneous singletons. Within this model, graph and coefficient
restrictions do not make exact feasibility easier than general
existential-real decision. The solution-preserving electrical realization
is the main contribution; the arithmetic and topological tools on which it
builds are established.

The same correspondence explains the algebraic and numerical behavior of
feasible voltage sets. Rational equivalence realizes precisely the compact
basic closed sets over $\Q$, semialgebraic homeomorphism realizes all
compact semialgebraic topologies, and unique voltages can have arbitrarily
high algebraic degree. Separately, the ordinary bounded-degree, fixed-data
construction gives infeasible networks with doubly exponentially small
injection residuals, the worst-case scale for that class; no
residual-preserving planarization is asserted. Exact decision,
tolerance-based numerical evidence, and certificates for a gap promise
therefore require different guarantees.

For AC models, the transfer depends on the meaning of an angle limit. An
explicit rational quadratic formula encodes the established zero-winding
condition using linearly many variables and predicates. Reactive balance then gives an exact equal-angle
theorem and a quantitative stability estimate. Principal-angle constraints
alone allow winding states, whereas bus-angle boxes and size-dependent
principal windows support separate hardness transfers. These conclusions
classify exact feasibility for the specified models and delimit
residual-based certification. They do not predict typical solver
performance or cover operating models that exclude the prescribed voltage
and injection combinations used by the reduction.
